\section{Conclusiones}
La crisis territorial en Santa Cruz de Tenerife no viene impulsada por una expansión demográfica descontrolada, sino por un embotellamiento normativo artificial:
un marco urbanístico obsoleto y judicialmente paralizado,
incapaz de dar cabida incluso a la modesta formación de nuevos hogares, a la modernización económica y a las transiciones terciarias.

La situación actual del planeamiento urbanístico de Santa Cruz de Tenerife se encuentra jurídicamente fragmentada,
dado que diversas resoluciones judiciales han anulado sistemáticamente los planes generales más recientes, así como sus adaptaciones provisionales.

Esta inestabilidad en la planificación urbana genera una importante inseguridad jurídica, especialmente en lo que respecta a las licencias de edificación.
Sin un plan integral vigente y actualizado, el marco legal para la concesión de licencias se vuelve ambiguo.
Si bien algunas fuentes sugieren la aplicación parcial del marco normativo de 2005 o de regulaciones transitorias,
la anulación de la ordenanza de 2022 ha socavado, en particular, el intento más reciente de establecer normas claras para el uso del suelo. Esta situación exige extremar las precauciones ante cualquier proyecto urbanístico, ya que la validez de los permisos podría ser impugnada debido a la contradicción entre los estatus jurídicos de los distintos instrumentos de planificación.
